\usepackage{fontspec}
% La norma exige Times New Roman. No se distribuye una fuente propietaria.
% Opción 1: cuatro archivos propios en fuentes/ (ver README).
\IfFileExists{fuentes/times.ttf}{
 \setmainfont{times.ttf}[Path=fuentes/,BoldFont=timesbd.ttf,ItalicFont=timesi.ttf,BoldItalicFont=timesbi.ttf]
}{\IfFontExistsTF{Times New Roman}{\setmainfont{Times New Roman}}{
 \iffuenteestricta
  \PackageError{memoria}{Falta Times New Roman}{Suba los cuatro TTF propios a fuentes/.}
 \fi
 \IfFontExistsTF{TeX Gyre Termes}{\setmainfont{TeX Gyre Termes}}{\setmainfont{NimbusRoman-Regular.otf}[BoldFont=NimbusRoman-Bold.otf,ItalicFont=NimbusRoman-Italic.otf,BoldItalicFont=NimbusRoman-BoldItalic.otf]}
 \PackageWarningNoLine{memoria}{Fuente alternativa de previsualizacion. Instale Times New Roman antes de la entrega}
}}
\usepackage[english]{babel}
\babelprovide[main,import]{spanish}
\addto\captionsspanish{\renewcommand{\tablename}{Tabla}\renewcommand{\appendixname}{Anexo}}
\usepackage[left=4cm,right=3cm,top=3cm,bottom=3cm,footskip=1.2cm]{geometry}
\usepackage{setspace}
\onehalfspacing
% 12 pt: espacio natural 14.5 pt; 1.5 interlineado = 18.125 pt.
% Añadir 6.042 pt entre párrafos da 24.167 pt: doble interlineado.
\setlength{\parskip}{6.042pt plus 1pt minus 1pt}
\setlength{\parindent}{1.25cm}
\usepackage{amsmath,amssymb}
\usepackage{graphicx,xcolor,booktabs,array,tabularx,longtable}
\usepackage{caption}
\captionsetup{font=small,labelfont=bf}
\captionsetup[table]{position=top}
\usepackage{float,enumitem,titlesec,fancyhdr,etoolbox}
\usepackage{tikz}
\usetikzlibrary{positioning,arrows.meta}
\usepackage{xurl}
\usepackage[hidelinks,unicode]{hyperref}
\hypersetup{pdftitle={\TituloMemoria},pdfauthor={\AutorMemoria},pdfsubject={Memoria de titulación: evaluación de cuantización en microcontroladores}}
\pagestyle{fancy}\fancyhf{}\fancyfoot[R]{\thepage}
\renewcommand{\headrulewidth}{0pt}
\fancypagestyle{plain}{\fancyhf{}\fancyfoot[R]{\thepage}\renewcommand{\headrulewidth}{0pt}}
\titleformat{\chapter}[block]{\normalfont\bfseries\Large}{\thechapter}{.7em}{}
\titleformat{name=\chapter,numberless}[block]{\normalfont\bfseries\Large\filcenter}{}{0pt}{}
\titlespacing*{\chapter}{0pt}{0pt}{20pt}
\titleformat{\section}{\normalfont\bfseries\large}{\thesection}{.7em}{}
\titleformat{\subsection}{\normalfont\bfseries}{\thesubsection}{.7em}{}
% Un nivel adicional para conservar la jerarquía de los antiguos capítulos.
\setcounter{secnumdepth}{3}\setcounter{tocdepth}{3}
\titleformat{\subsubsection}{\normalfont\bfseries}{\thesubsubsection}{.7em}{}
\addto\captionsspanish{\renewcommand{\contentsname}{ÍNDICE}\renewcommand{\listfigurename}{ÍNDICE DE FIGURAS}\renewcommand{\listtablename}{ÍNDICE DE TABLAS}\renewcommand{\bibname}{REFERENCIAS}}
\setlist{itemsep=2pt,topsep=4pt}
\renewcommand{\arraystretch}{1.16}
\newcolumntype{Y}{>{\raggedright\arraybackslash}X}
\newcommand{\preliminar}[1]{\chapter*{\centering #1}\phantomsection\addcontentsline{toc}{chapter}{#1}}
\newcommand{\pendiente}[1]{\ifnotaseditoriales\par\noindent{\small\color{brown}\textbf{POR COMPLETAR:} #1}\par\fi}
\newcommand{\notaeditorial}[1]{\ifnotaseditoriales\par\noindent{\small\color{blue!55!black}\textbf{NOTA DE EDICIÓN:} #1}\par\fi}
\newcommand{\ND}{\textemdash}
\newcommand{\termino}[1]{\textquotedblleft#1\textquotedblright}
\newcommand{\fuente}[1]{\par\vspace{1pt}{\footnotesize Fuente: #1.}\par}
\newcommand{\figurapendiente}[3]{\begin{figure}[H]\centering
\fbox{\parbox[c][3.2cm][c]{.90\linewidth}{\centering\small Figura por incorporar\\[5pt]#1}}
\caption{#2}\label{#3}\end{figure}}
% Bibliografía sencilla, numérica y en orden de citación. Revisar campos según anexo C.
\AtBeginEnvironment{quote}{\singlespacing} % etoolbox se carga explícitamente abajo.
\raggedbottom
\emergencystretch=2em

% Figuras del capítulo Resultados: tamaño legible y altura limitada.
\newcommand{\resfig}[5]{\begin{figure}[htbp]\centering
\includegraphics[width=#5\linewidth,height=.74\textheight,keepaspectratio]{figuras/resultados/#1}
\caption{#2}\label{#3}\fuente{#4}\end{figure}}

\usepackage{placeins}
